\chapter{Εισαγωγή}\label{ch:intro}

\section{Παρουσίαση Θέματος}
Η επίσημη περιγραφή της πτυχιακής αναφέρει: {\itshape Η πτυχιακή εργασία ασχολείται με την ανάπτυξη μιας ασφαλούς και αποδοτικής \ENG{SCORM Engine},
η οποία θα υποστηρίζει τα πιο διαδεδομένα \ENG{SCORM types} και θα ενσωματώνει μηχανισμούς αποθήκευσης και συγχρονισμού της προόδου των χρηστών.
Το σύστημα θα επικοινωνεί με \ENG{Learning Management Systems (LMS)} για τη μετάδοση των δεδομένων προόδου, ενώ η υλοποίηση δίνει έμφαση στην ασφάλεια,
την επεκτασιμότητα και την αποδοτική διαχείριση των εκπαιδευτικών δεδομένων}. 

Με αφορμή τη δημιουργία μιας τέτοιας μηχανής, πραγματοποιείται εκτενής ανάλυση του θεωρητικού και ιστορικού πλαισίου, ώστε να διαμορφωθεί στέρεο υπόβαθρο τόσο 
για τον συγγραφέα όσο και για οποιονδήποτε επιχειρήσει να επεκτείνει το παρόν έργο. Παράλληλα, το κείμενο λειτουργεί ως συμπαγής βιβλιοθήκη γνώσεων για το αντικείμενο.
Επιπλέον, ένας ακόμη λόγος για την εκπόνηση της πτυχιακής είναι η διεύρυνση των γνώσεων του συγγραφέα στο συγκεκριμένο πεδίο.

Οι πληροφορίες αντλήθηκαν κυρίως από επίσημες, πηγές προδιαγραφές και τεκμηριώσεις (\ENG{Documentation}), συμπληρωμένες από επαγγελματική εμπειρία του συγγραφέα στο αντικείμενο.

\section{Σκοπός \& Στόχοι}
Κύριος σκοπός της διπλωματικής εργαίας είναι εκτέλεση \ENG{SCORM} πακέτων και αξιόπιστη επικοινωνία προόδου με οποιοδήποτε \ENG{LMS}.

\noindent \textbf{Επιμέρους στόχοι:}
\begin{itemize}
  \item \textbf{Συμβατότητα \& διαλειτουργικότητα:} Υποστήριξη των πιο διαδεδομένων εκδόσεων \ENG{SCORM} (\ENG{1.2} και \ENG{2004}), με σωστή υλοποίηση \ENG{API Runtime} και \ENG{Data Model}.
  \item \textbf{Ασφάλεια περιεχομένου:} Απομόνωση εκτέλεσης (\ENG{sandboxing}), έλεγχοι έναντι \ENG{XSS}/\ENG{CSRF} και πολιτικές περιεχομένου (\ENG{CSP}) όπου αρμόζει.
  \item \textbf{Αποθήκευση \& συγχρονισμός:} Αξιόπιστη επιμονή δεδομένων προόδου (\ENG{persistence}) και ανθεκτικός συγχρονισμός με \ENG{LMS}, ακόμη και σε συνθήκες διακοπτόμενης συνδεσιμότητας.
  \item \textbf{Επεκτασιμότητα \& παρατηρησιμότητα:} Σχεδίαση για κλιμάκωση, με μετρικές, καταγραφή συμβάντων και δυνατότητες παρακολούθησης/εντοπισμού σφαλμάτων (επαρκές και κατανοητό \ENG{Logging}).
  \item \textbf{Αναφορές \& \ENG{analytics}:} Παροχή βασικών δεικτών προόδου/ολοκλήρωσης και διεπαφών εξαγωγής δεδομένων για περαιτέρω ανάλυση.
  \item \textbf{Τεκμηρίωση \& επαναχρησιμοποίηση:} Καθαρή τεκμηρίωση αρχιτεκτονικής και κώδικα, ώστε το έργο να επεκτείνεται εύκολα από τρίτους.
\end{itemize}

\section{Σημασία}
Η ανάπτυξη μιας σύγχρονης \ENG{SCORM Engine} παραμένει ουσιαστική για τον χώρο της ηλεκτρονικής μάθησης (\ENG{E-Learning}):
\begin{itemize}
  \item \textbf{Μεγάλη εγκατεστημένη βάση περιεχομένου:} Σημαντικό εκπαιδευτικό υλικό παραμένει σε \ENG{SCORM} (ή άλλα \ENG{e-learning standards}), καθιστώντας κρίσιμη τη σωστή εκτέλεση και παρακολούθηση προόδου.
  \item \textbf{Διαλειτουργικότητα:} Εξασφαλίζει ότι περιεχόμενο και δεδομένα κινούνται ομαλά μεταξύ διαφορετικών \ENG{LMS} και υποδομών.
  \item \textbf{Ασφάλεια σε ετερογενή περιβάλλοντα:} Η ασφαλής εκτέλεση «μη έμπιστου» εκπαιδευτικού περιεχομένου μειώνει τον κίνδυνο επιθέσεων και διαρροών.
  \item \textbf{Κλιμάκωση \& αξιοπιστία:} Η υποστήριξη μεγάλων φορτίων χρηστών απαιτεί σωστό σχεδιασμό επιμονής, ουρών/επανάληψης και παρατηρησιμότητας.
  \item \textbf{Αξιοποίηση δεδομένων μάθησης:} Συγκεντρωτικά στοιχεία (\ENG{analytics}) ενισχύουν τον παιδαγωγικό σχεδιασμό και τη βελτίωση μαθησιακών αποτελεσμάτων.
  \item \textbf{Προσωπική \& επαγγελματική εξέλιξη:} Η υλοποίηση επιτρέπει στον γράφοντα να εμβαθύνει στο αντικείμενο, να εφαρμόσει βέλτιστες πρακτικές λογισμικού και να επεκτείνει ουσιαστικά τις γνώσεις του.
\end{itemize}

\section{Μεθοδολογία}
\todoENG{Σύντομη περιγραφή.}

\section{Περιορισμοί \& Υποθέσεις}
\begin{itemize}
  \item Εστίαση σε \ENG{SCORM 1.2} και \ENG{2004}. Μελλοντικά μπορεί να γίνει επέκταση για περισσοτερα \ENG{e-learning standards}.
  \item Θεωρούμε ότι το \ENG{engine} θα λειτουργεί ως ξεχωριστή εφαρμογή 
\end{itemize}

\section{Δομή του Εγγράφου}
Το \Cref{ch:intro} περιγράφει εν συντομία εισαγωγικά στοιχεία που αφορούν τη συγκεριμένη διπλωματική εργασία.
Το \Cref{ch:history} παρουσιάζει το ιστορικό υπόβαθρο του \ENG{SCORM} και γενικότερα αυτής της μορφής μάθησης.